% Generated draft: requires mathematical and source review.
\documentclass[11pt]{article}
\usepackage[margin=25mm]{geometry}
\usepackage{fontspec}
\setmainfont{Latin Modern Roman}
\usepackage{amsmath,amssymb,mathtools,unicode-math}
\setmathfont{Latin Modern Math}
\usepackage{enumitem,needspace,xurl,hyperref}
\setlength{\parindent}{0pt}
\setlength{\parskip}{5pt}
\setlength{\emergencystretch}{3em}
\newcommand{\N}{\mathbb{N}}
\newcommand{\Z}{\mathbb{Z}}
\newcommand{\Q}{\mathbb{Q}}
\newcommand{\R}{\mathbb{R}}
\newcommand{\C}{\mathbb{C}}
\newcommand{\F}{\mathbb{F}}
\newcommand{\A}{\mathbb{A}}
\newcommand{\PP}{\mathbb{P}}
\newcommand{\E}{\mathbb{E}}
\newcommand{\eps}{\varepsilon}
\newcommand{\dd}{\,\mathrm{d}}
\begin{document}
\section*{UnsolvedMath ID: 30002793}
\textbf{Source problem code:} catalogue entry
\subsection{Definitions and mathematical statement}
The corrected target supplied for this record is the classification of arithmetically Cohen--Macaulay configurations of lines in projective three-space. The source does not specify a ground field, its characteristic, or whether it is algebraically closed. Accordingly, the formulation below is relative to a field \(k\); the intended restrictions on \(k\) remain unresolved.

Under the interpretation of a line configuration as a nonempty finite reduced union of distinct projective lines, let
\[
S=k[x_0,x_1,x_2,x_3],\qquad
X=\bigcup_{i=1}^{r}L_i\subseteq\mathbb P^3_k,\qquad r\geq 1,
\]
where each \(L_i\) is a projective linear subspace of dimension one defined over \(k\). Give \(X\) its reduced scheme structure, and set
\[
I(X)=\bigcap_{i=1}^{r}I(L_i),\qquad A_X=S/I(X).
\]
The configuration is \emph{arithmetically Cohen--Macaulay} (ACM) if its homogeneous coordinate ring \(A_X\) is Cohen--Macaulay. Since \(\dim A_X=2\), this is equivalent to
\[
\operatorname{depth}_{S_+}A_X=2,
\qquad S_+=(x_0,x_1,x_2,x_3).
\]
\textbf{Classification problem.} Characterize all such configurations \(\{L_1,\ldots,L_r\}\), for arbitrary positive integers \(r\), for which \(A_X\) is Cohen--Macaulay. The finite reduced-union interpretation and the intended ground-field hypotheses require confirmation from the source; no additional restrictions on intersections or positions of the lines are supplied.

\textbf{Source qualification.} The original registry wording also asked which reduced, possibly irreducible, curves \(C\subseteq\mathbb P^3\) satisfy
\[
\alpha\bigl(I(C)^{(2)}\bigr)-\alpha\bigl(I(C)\bigr)=1.
\]
Here \(\alpha(J)=\min\{d\geq 0:J_d\neq 0\}\) is the initial degree of a nonzero homogeneous ideal, and, for a radical homogeneous ideal \(I\), its second symbolic power is
\[
I^{(2)}=\bigcap_{\mathfrak p\in\operatorname{Min}(I)}\bigl(I^2S_{\mathfrak p}\cap S\bigr).
\]
The supplied statement-verification note assigns that curve question to the separate source code \texttt{OWR-13494-013}; it is not an additional condition on the configurations in this record. The supplied literature assessment discusses an ACM subcase and the unresolved existence of a non-ACM line configuration with this initial-degree gap. Its applicability to the unrestricted ACM classification problem above is unclear.
\subsection{Short English statement}
Characterize the finite reduced unions of lines in \(\mathbb P^3\) whose homogeneous coordinate rings are Cohen--Macaulay. The intended ground field and precise convention for a line configuration need confirmation. The supplied correction separates the symbolic-initial-degree question for reduced curves into \texttt{OWR-13494-013}.
\subsection{Sources}
\begin{itemize}
\item Local UnsolvedMath registry, supplied complete record with permanent ID \(30002793\), source problem code \texttt{catalogue entry}, in the Oberwolfach Reports Open Problems collection. The supplied proposer field is empty; the proposed year is 2015.
\item \emph{Mini-Workshop: Ideals of Linear Subspaces, Their Symbolic Powers and Waring Problems} (2015), Oberwolfach Reports, pp.~489--532, DOI \texttt{10.4171/owr/2015/9}. Supplied report and literature-source link: \url{https://doi.org/10.4171/owr/2015/9}.
\item Janssen, \emph{On the fattening of lines in P3}. Supplied literature-source link: \url{https://doi.org/10.1016/j.jpaa.2014.05.033}. No theorem number or theorem statement is supplied.
\item The supplied literature assessment is dated \texttt{2026-08-21}. It labels the problem open and states that an ACM subcase is classified by a cited theorem, while the existence of a non-ACM line configuration with the stated initial-degree gap remains unresolved. This is supplied dated curation derived from the OWR material, not an independently verified literature assessment; its relation to the corrected classification target requires review.
\item The registry labels the statement \texttt{corrected\_verified} and reports that the official report separates the ACM classification question from the curve symbolic-initial-degree question, assigning the latter to \texttt{OWR-13494-013}. These are supplied verification claims; the linked sources have not been independently consulted for this draft.
\end{itemize}
\end{document}
